\chapter{Test e Validazione}

In questo capitolo vengono presentati i risultati delle sessioni di monitoraggio effettuate sull'emulatore Android Cuttlefish. Per ciascuno scenario, le sette probe del sistema sono state eseguite in parallelo: la probe delle syscall per il monitoraggio di \texttt{openat}, \texttt{connect} ed \texttt{execve}; la probe del Binder per le transazioni IPC; la probe del ciclo di vita dei processi per la nascita e terminazione; la probe di associazione socket e la probe di accettazione connessioni per il monitoraggio dei socket in ascolto e delle connessioni in ingresso; la probe di invio dati e la probe di ricezione dati per il volume di dati inviati e ricevuti.

Gli scenari analizzati sono tre: l'attivazione della modalità aereo, una sessione di navigazione web con Chrome, e l'apertura dell'applicazione galleria con accesso alla fotocamera.

% -------------------------------------------------------
\section{Attivazione della modalità aereo}
\label{sec:aereo}

L'applicazione monitorata è \texttt{com.android.settings} (uid 1000), il pannello di configurazione di sistema di Android. Lo scenario consiste nell'attivare e poi disattivare la modalità aereo tramite il menu delle impostazioni rapide. Il riferimento temporale è l'avvio della probe del Binder, con le altre sei probe avviate contestualmente. I due eventi chiave sono l'attivazione della modalità aereo e la sua disattivazione.

\begin{figure}[h]
    \centering
    \includegraphics[width=\textwidth]{images/timeline-aereo.png}
    \caption{Sequenza delle azioni nella sessione di attivazione e disattivazione della modalità aereo.}
    \label{fig:timeline-aereo}
\end{figure}

\subsection{Fase preliminare}

Nei 100 secondi precedenti all'attivazione, il sistema opera nella schermata home con il WiFi attivo. La probe del Binder registra un'attività di fondo dominata dal canale chiamato da \texttt{surfaceflinger}, con 3.330 chiamate per circa 1.4~MB: è il compositor grafico~\cite{surfaceflinger_docs} che aggiorna continuamente la schermata.

\subsection{Attivazione della modalità aereo}

All'attivazione della modalità aereo si produce un effetto immediato e netto sullo stack di rete. La probe di associazione socket non registra alcun evento tra durante l'intera finestra che copre il periodo in modalità aereo e la successiva riattivazione non genera alcun \texttt{bind} da processi di rete. È la conferma diretta che lo stack WiFi è stato completamente fermato.

Il percorso del comando attraverso il sistema è ricostruibile dalle probe. L'utente interagisce con \texttt{com.android.settings}, che invia via Binder~\cite{binder_ipc} il \textit{broadcast} di sistema \texttt{ACTION\_AIRPLANE\_MODE\_CHANGED} al \texttt{ConnectivityService} all'interno del \textit{system server}. È quest'ultimo, non l'applicazione impostazioni, a coordinare la sequenza di disattivazione: WiFi e Bluetooth vengono fermati attraverso chiamate ai rispettivi HAL via Binder. Il processo \texttt{ndroid.settings} accumula 4.169 transazioni Binder nell'intera sessione, ma nessuna di esse produce direttamente traffico di rete: tutto passa attraverso il \textit{system server}.

\subsection{Disattivazione e riconnessione}

Alla disattivazione della modalità aereo viene avviata una sequenza di riconnessione osservabile passo per passo attraverso le probe. Il primo segnale misurabile arriva con un po' di ritardo rispetto al toggle: vengono registrate le prime operazioni \texttt{bind} dopo la pausa. Successivamente \texttt{IpClient.wlan0} apre sei nuovi socket per avviare la negoziazione DHCP, e \texttt{netd} crea un thread per la risoluzione dei nomi. Il ritardo riflette il tempo necessario al \textit{firmware} dell'interfaccia WiFi virtuale di Cuttlefish per completare la riassociazione alla rete prima che lo stack di rete superiore possa riattivarsi.

La probe di ricezione dati mostra che \texttt{IpClient.wlan0} riceve complessivamente circa 229~KB nell'intera sessione: è il traffico DHCP~\cite{dhcp_rfc} corrispondente all'acquisizione dell'indirizzo IP dopo la riconnessione.

Il picco di Binder dell'intera sessione (4.099 eventi/s) si registra due secondi prima del primo bind di rete. Esso è prodotto dai \textit{broadcast} di \texttt{ConnectivityService} che notificano tutti i processi del sistema del cambiamento di stato della rete, prima ancora che il collegamento fisico sia ristabilito. Anche questa sequenza, il Binder che esplode prima che i socket di rete si riaprino, è osservabile solo grazie al monitoraggio parallelo delle sette probe.

\subsection{Osservazioni sull'architettura osservata}
\label{subsec:aereo-arch}

Lo scenario della modalità aereo si distingue rispetto ai prossimi due per la natura dell'evento monitorato: non si tratta dell'avvio di un'applicazione ma di un cambiamento di stato del sistema che coinvolge l'infrastruttura di rete.

Il segnale più caratteristico è la finestra di silenzio nelle probe di rete. La probe di associazione socket non registra eventi, e la probe di accettazione connessioni mostra un gap analogo. Questa assenza prolungata di attività, in un sistema che normalmente produce bind e accept continuamente, è l'indicatore diretto della disattivazione dello stack WiFi.

Il traffico inviato via socket (2.8~MB totali) è dominato non dal WiFi ma dalle comunicazioni interne dell'hardware audio e grafico (\texttt{android.hardwar}: 1.8~MB, \texttt{app} SurfaceControl: 709~KB). Il contributo diretto di \texttt{wificond} è di soli 4.9~KB, segno che il controllo della radio WiFi avviene quasi interamente via Binder e non via socket.

\begin{table}[h]
\centering
\begin{tabular}{lr}
\toprule
\textbf{Categoria} & \textbf{Volume totale} \\
\midrule
Transazioni Binder & 35.058 \\
Traffico IPC & 11.3 MB \\
Dati inviati via socket & 2.8 MB \\
Dati ricevuti via socket & 4.8 MB \\
Bind durante modalità aereo & 0 \\
\bottomrule
\end{tabular}
\caption{Riepilogo del traffico complessivo della sessione modalità aereo: il silenzio nelle probe di rete durante la finestra t$\,\approx\,$92s--147s è il segnale diretto della disattivazione dello stack WiFi.}
\end{table}

% -------------------------------------------------------
\section{Navigazione web con Chrome}
\label{sec:chrome}

L'applicazione monitorata è \texttt{org.chromium.webview\_shell} (uid 10064), un'applicazione di test che espone direttamente il motore di rendering WebView di Chromium. Lo scenario consiste nella navigazione web a partire dalla schermata home dell'emulatore: viene prima effettuata una ricerca tramite il widget di ricerca presente nella home, poi si naviga sul sito google.com e si esegue una ricerca di notizie. Le probe vengono avviate prima di qualsiasi interazione con il browser, per poter documentare il comportamento di base del sistema e distinguerlo dall'attività generata dall'utente.

La sessione si articola in quattro fasi, identificate in base alle azioni effettuate:

\begin{figure}[h]
    \centering
    \includegraphics[width=\textwidth]{images/timeline-rete.png}
    \caption{Sequenza delle azioni nella sessione di navigazione web.}
    \label{fig:timeline-chrome}
\end{figure}

\subsection{Fase preliminare}

Nei primi 36 secondi l'emulatore si trova nella schermata home e non viene effettuata alcuna interazione. Questa fase serve a documentare il comportamento di base del sistema: Android non è silenzioso nemmeno in \textit{idle}, e senza questa \textit{baseline} sarebbe impossibile distinguere il traffico di fondo dall'attività riconducibile alle azioni utente.

La probe di ricezione dati mostra traffico già dai primi secondi, generato da \texttt{logcat} e da altri servizi di sistema che comunicano su socket locali in modo continuativo. Più interessante è un picco di invio che compare intorno a t$\,\approx\,$27s, con il sistema che trasmette prima 6.4~KB poi 28.9~KB in due secondi consecutivi. Il traffico è dominato da \texttt{InputDispatcher} (18.4~KB a t$\,\approx\,$28s) e dal processo \texttt{app} (SurfaceControl), con un contributo minore di \texttt{android.ui}; non ha nulla a che fare con il browser, che non è ancora avviato. Si tratta di comunicazioni interne al sottosistema grafico e di input, un esempio di quanto il sistema operativo produca traffico di rete per ragioni completamente indipendenti dall'utente.

\begin{figure}[h]
    \centering
    \includegraphics[width=\textwidth]{images/picco-invii.png}
    \caption{Picco di invio nella fase preliminare, generato dal compositor grafico prima di qualsiasi interazione utente.}
    \label{fig:picco-invii}
\end{figure}

\subsection{Invio della query}

Il click sul widget di ricerca non produce eventi netti nelle probe di rete. L'attività misurabile parte qualche secondo dopo, quando la probe delle syscall intercetta una \texttt{connect} da parte dell'applicazione Google Search che gestisce il widget nella schermata home. Contestualmente, la probe di accettazione connessioni mostra che \texttt{netd} crea un thread che  nasce per gestire la sua richiesta DNS. È quindi l'app di ricerca Google, e non il browser, ad avviare la risoluzione del nome mentre l'utente digita la query.

La prima attività di rete del processo Chrome compare circa nove secondi dopo il click, quando le probe di invio e ricezione dati registrano i primi eventi dal thread \texttt{RenderThread} (PID 8373). I thread \texttt{m.webview\_shell} e \texttt{NetworkService} compaiono nei secondi successivi. Il fatto che thread con nomi diversi condividano lo stesso PID è un effetto del limite di 15 caratteri del campo \texttt{comm} nel kernel Linux~\cite{linux_task_struct}, che riporta il nome del singolo thread piuttosto che quello del processo.

Successivamente, la probe del Binder registra i primi eventi associati a \url{CrRenderer\_Main} (PID 8406),
il renderer di Chromium. La comunicazione tra \texttt{NetworkService}
e il renderer avviene via Binder, non via socket: nell'arco dell'intera sessione vengono registrate 1031 transazioni tra i 
due per un totale di circa 339~KB. Ogni risposta HTTP ricevuta dalla rete viene consegnata al renderer attraverso questo 
canale IPC prima di essere visualizzata.

\subsection{Caricamento di google.com}

Il click su google.com produce in un secondo un \textit{burst} di otto operazioni \texttt{bind} da parte di \texttt{NetworkService}, che indicano l'apertura di nuovi socket per il caricamento parallelo delle risorse della pagina, coerente con il \textit{multiplexing} HTTP/2~\cite{http2_rfc} utilizzato da Chrome.

Il picco di connessioni in ingresso arriva con un ritardo di circa dodici secondi rispetto al click. Il ritardo è compatibile con il tempo di caricamento del contenuto dinamico della homepage di Google, che include numerose richieste secondarie per immagini, script e dati personalizzati.

\subsection{Ricerca ``google news''}

La ricerca di ``google news'' produce la risposta più netta dell'intera sessione. La probe delle syscall registra il picco assoluto di 608 eventi al secondo, contro una media di circa 187 per l'intera sessione.

Il picco di invio dati segue con un ritardo di circa dodici secondi: la probe di invio dati registra 79.4~KB inviati in un secondo, il valore massimo dell'intera sessione. Il ritardo riflette il tempo che il browser impiega a ricevere i risultati, elaborarli tramite il renderer e inviare le richieste secondarie per i contenuti collegati.

\begin{figure}[h]
    \centering
    \includegraphics[width=\textwidth]{images/send-peak.png}
    \caption{Picco di 79.4~KB inviati in un secondo, valore massimo dell'intera sessione.}
    \label{fig:send-peak}
\end{figure}

\subsection{Osservazioni sull'architettura osservata}
\label{subsec:chrome-arch}

Il monitoraggio parallelo delle sette probe ha permesso di osservare una catena di eventi che nessuna singola probe avrebbe potuto documentare per intero.

Il renderer \texttt{CrRendererMain} (PID 8406) non ha accesso diretto alla rete. Tutta la comunicazione tra il codice di rendering e lo stack di rete passa attraverso Binder: la catena osservata è \texttt{m.webview\_shell} $\rightarrow$ Binder $\rightarrow$ \texttt{NetworkService} $\rightarrow$ socket. Questo isolamento è una scelta progettuale di Chromium~\cite{chromium_multiprocess} per separare il processo di rendering, potenzialmente esposto a contenuti non fidati, dal processo che gestisce le connessioni di rete.

Complessivamente, la sessione ha prodotto 42\,343 transazioni Binder per un totale di circa 17~MB trasferiti via IPC, a fronte di 1.55~MB inviati e 2.37~MB ricevuti via socket. Il volume di traffico IPC interno supera quindi di circa cinque volte il traffico di rete esterno: anche una semplice sessione di navigazione web genera internamente un volume di comunicazione inter-processo di gran lunga superiore a quello scambiato con la rete.

\begin{table}[h]
\centering
\begin{tabular}{lr}
\toprule
\textbf{Categoria} & \textbf{Volume totale} \\
\midrule
Transazioni Binder & 42.343 \\
Traffico IPC & 17 MB \\
Dati inviati via socket & 1.55 MB \\
Dati ricevuti via socket & 2.37 MB \\
\bottomrule
\end{tabular}
\caption{Riepilogo del traffico complessivo della sessione Chrome: 42\,343 transazioni Binder e confronto tra volume IPC e traffico di rete.}
\end{table}

% -------------------------------------------------------
\section{Apertura della galleria e scatto fotografico}
\label{sec:galleria}

L'applicazione monitorata è l'applicazione predefinita per la gestione delle immagini su Android (uid 10066). Lo scenario combina due applicazioni distinte: nella prima parte l'utente naviga la galleria fotografica, nella seconda passa alla fotocamera e scatta un'immagine. Il riferimento temporale è l'avvio della probe del Binder, con le altre sei probe avviate contestualmente nei secondi precedenti. Gli eventi si distribuiscono in tre momenti precisi: apertura della galleria, attivazione della fotocamera e scatto della foto.

\begin{figure}[h]
    \centering
    \includegraphics[width=\textwidth]{images/timeline-galleria.png}
    \caption{Sequenza delle azioni nella sessione di apertura galleria e fotocamera.}
    \label{fig:timeline-galleria}
\end{figure}

\subsection{Apertura della galleria}

All'apertura dell'applicazione galleria, la probe del ciclo di vita dei processi registra la creazione del processo \texttt{droid.gallery3d} (PID 14805) come figlio del \textit{system server} di Android. L'associazione all'uid 10066 è confermata dai percorsi dei file \texttt{cgroup} che il processo \texttt{ActivityManager} monitora per la gestione della memoria dell'applicazione.

La probe delle syscall rileva le prime operazioni di accesso allo storage: un thread dell'applicazione apre la radice dello spazio di archiviazione condiviso. Questo accesso corrisponde alla scansione iniziale del catalogo di immagini che la galleria esegue all'avvio per costruire la lista dei contenuti disponibili.

Sul fronte del Binder, l'avvio della galleria produce un'attività crescente che culmina nel picco assoluto della sessione: la probe registra 4\,723 transazioni in un singolo secondo, circa cinque volte la media della sessione (941 eventi/s). Questo picco corrisponde al momento in cui la galleria carica le anteprime delle immagini, operazione che richiede un'intensa sequenza di comunicazioni con \texttt{SurfaceFlinger}~\cite{surfaceflinger_docs} per il rendering di ciascun thumbnail.

\begin{figure}[h]
    \centering
    \includegraphics[width=\textwidth]{images/binder-galleria-peak.png}
    \caption{Picco di 4\,723 transazioni Binder, durante il caricamento delle anteprime da parte della galleria.}
    \label{fig:binder-galleria-peak}
\end{figure}

\subsection{Attivazione della fotocamera}

Il passaggio alla fotocamera produce il cambiamento più caratteristico dell'intera sessione nel profilo delle comunicazioni Binder. Il processo \texttt{android.camera2} (PID 1913) diventa attivo e due thread legati all'hardware della fotocamera emergono tra i processi più trafficati: \texttt{Cam0\_ResultDisp} con 6\,285 eventi e \texttt{C3Dev-0-ReqQueu} con 4\,280 eventi, il secondo e il quarto posto nell'intera sessione.

Il grafo IPC mostra che \texttt{Cam0\_ResultDisp} comunica principalmente con il canale che corrisponde alla pipeline di elaborazione dei frame della fotocamera: il thread che riceve i risultati dal sensore li trasmette al \texttt{Camera Hardware Abstraction Layer}~\cite{camera_hal} attraverso Binder, che li consegna al processo di visualizzazione per l'anteprima in tempo reale. L'attivazione di questo canale è assente in tutti gli altri scenari analizzati e rappresenta quindi un segnale affidabile dell'accesso all'hardware della fotocamera.

In parallelo, \texttt{SurfaceFlinger} mantiene il canale più trafficato dell'intera sessione: 4\,481 chiamate verso il \texttt{binder} per un totale di circa 2.1~MB. Questo traffico riflette la composizione continua dei frame dell'anteprima, che richiede una frequenza di comunicazione elevata tra il compositor grafico e il driver del display.

\begin{figure}[h]
    \centering
    \includegraphics[width=\textwidth]{images/binder-camera-ipc.png}
    \caption{Grafo delle comunicazioni Binder durante l'anteprima della fotocamera: il canale \texttt{Cam0\_ResultDisp} verso il Camera HAL è il segnale caratteristico dell'accesso all'hardware fotografico.}
    \label{fig:binder-camera-ipc}
\end{figure}

\subsection{Scatto della foto}

Lo scatto della fotografia  lascia una traccia diretta nella resource map prodotta dalla probe delle syscall. Un thread del processo \texttt{droid.gallery3d} apre il file \texttt{IMG\_20260304\_113342.jpg}: si tratta dell'immagine appena acquisita, letta dall'applicazione per aggiornare il catalogo con il nuovo scatto. Lo stesso thread accede anche a \texttt{/share\_history.xml}, il file che l'applicazione aggiorna sistematicamente a ogni modifica del catalogo.

La probe di ricezione dati registra circa 215~KB ricevuti dal processo \texttt{droid.gallery3d}  su socket locali. Il ritardo di circa tredici secondi rispetto allo scatto è compatibile con il tempo necessario al sistema di gestione dei media di Android per processare il nuovo file, generare l'anteprima e notificare l'applicazione dell'avvenuto aggiornamento del catalogo.

Nessuna comunicazione di rete esterna è associata allo scatto: la probe di invio dati registra per \texttt{droid.gallery3d} un totale di soli circa 4~KB durante l'intera sessione, esclusivamente su socket locali verso servizi di sistema.

\subsection{Osservazioni sull'architettura osservata}
\label{subsec:galleria-arch}

Lo scenario galleria e fotocamera si distingue dalla navigazione web per l'assenza totale di traffico di rete esterno. La probe di associazione socket non registra alcun evento per l'intera sessione, e quella di accettazione connessioni rileva attività solo da processi di sistema come \texttt{init} e \texttt{netd}, mai dalle applicazioni utente. È la conferma che galleria e fotocamera operano esclusivamente su risorse locali.

Nonostante questa chiusura alla rete, il volume di comunicazione IPC è considerevole: la sessione ha prodotto 38\,824 transazioni Binder per un totale di circa 17.2~MB trasferiti via IPC, a fronte di appena 4~KB inviati via socket. Anche operazioni che l'utente percepisce come completamente locali generano internamente un flusso continuo di messaggi tra decine di componenti di sistema.

Il dato più interessante riguarda l'identificazione dell'accesso alla fotocamera. L'attivazione del thread \texttt{Cam0\_ResultDisp} e del canale Binder verso il Camera HAL è un indicatore preciso e stabile: non compare durante l'uso della galleria, appare non appena la fotocamera viene avviata, e persiste fino alla chiusura dell'applicazione. Questo tipo di segnatura IPC potrebbe essere utilizzato per rilevare l'accesso alla fotocamera da parte di un processo, indipendentemente da qualsiasi indicatore a livello applicativo.

\begin{table}[h]
\centering
\begin{tabular}{lr}
\toprule
\textbf{Categoria} & \textbf{Volume totale} \\
\midrule
Transazioni Binder & 38\,824 \\
Traffico IPC & 17.2 MB \\
Dati inviati via socket & 4 KB \\
Dati ricevuti via socket & 215 KB \\
Bind aperti & 0 \\
\bottomrule
\end{tabular}
\caption{Riepilogo del traffico complessivo della sessione galleria e fotocamera: 38\,824 transazioni Binder e traffico di rete trascurabile.}
\end{table}
